% Generated draft: requires mathematical and source review.
\documentclass[11pt]{article}
\usepackage[margin=25mm]{geometry}
\usepackage{fontspec}
\setmainfont{Latin Modern Roman}
\usepackage{amsmath,amssymb,mathtools,unicode-math}
\setmathfont{Latin Modern Math}
\usepackage{enumitem,needspace,xurl,hyperref}
\setlength{\parindent}{0pt}
\setlength{\parskip}{5pt}
\setlength{\emergencystretch}{3em}
\newcommand{\N}{\mathbb{N}}
\newcommand{\Z}{\mathbb{Z}}
\newcommand{\Q}{\mathbb{Q}}
\newcommand{\R}{\mathbb{R}}
\newcommand{\C}{\mathbb{C}}
\newcommand{\F}{\mathbb{F}}
\newcommand{\A}{\mathbb{A}}
\newcommand{\PP}{\mathbb{P}}
\newcommand{\E}{\mathbb{E}}
\newcommand{\eps}{\varepsilon}
\newcommand{\dd}{\,\mathrm{d}}
\begin{document}
\section*{UnsolvedMath ID: 30000886}
\textbf{Source problem code:} catalogue entry
\subsection{Definitions and mathematical statement}
For a finite non-Hamiltonian group \(G\), the first question defines \(m(G)\) to be the least positive integer \(m\) such that, for every pair of bicyclic units \(b_1,b_2\) with \(\langle b_1,b_2\rangle\) not nilpotent, the subgroup \(\langle b_1,b_2^m\rangle\) is free. It asks whether there is an absolute integer \(M\) such that \(m(G)\leq M\) for every such finite group. Existence of the minimum for each group is implicit in the source and requires its theorem context.

In the customary integral group-ring setting, bicyclic units have forms \(1+(1-a)b\widehat a\) or \(1+\widehat a b(1-a)\), where \(a,b\in G\), \(a\) has finite order \(r\), and \(\widehat a=1+a+\cdots+a^{r-1}\). The source does not specify the coefficient ring or whether both types are included. ``Hamiltonian'' usually means nonabelian with every subgroup normal; its exact convention here should also be checked. The intended freeness is normally nonabelian freeness on the indicated generators, not merely an unspecified abstract cyclic free group.

The record calls the following related algebraic-integer question equivalent. For a nonzero algebraic integer \(z\in\C\), it describes a free point using the matrices
\[
U=\begin{pmatrix}1&1\\0&1\end{pmatrix},\qquad
V(z)=\begin{pmatrix}1&z\\0&1\end{pmatrix},
\]
and defines \(m(z)\) as the least positive integer \(k\) for which \(kz\) is a free point. It asks whether these minima are uniformly bounded over all nonzero algebraic integers.

There is a material defect in this supplied matrix formulation: both displayed matrices are upper-unitriangular and commute. Thus they cannot freely generate a nonabelian group of rank two, and the stated predicate cannot define the intended nonabelian free points. Even the meaning of ``free'' and the existence of \(m(z)\) are therefore unresolved as written. A different off-diagonal placement might have been intended, but no matrix entry is silently changed. The claimed equivalence with the bicyclic-unit question is recorded as a source claim and cannot be used until the matrices and reduction are checked.
\subsection{Short English statement}
Are the powers needed to obtain free subgroups from non-nilpotent pairs of bicyclic units uniformly bounded? The related algebraic-integer formulation is defective as supplied: its two unipotent matrices commute, so they cannot generate a rank-two free group.
\subsection{Sources}
\begin{itemize}
\item Local UnsolvedMath registry, permanent ID \(30000886\), source code catalogue entry, ``Uniform Freeness Bounds for Bicyclic Units.'' No individual proposer is supplied.
\item \emph{Mini-Workshop: Arithmetik von Gruppenringen} (2008), Oberwolfach Reports, DOI 10.4171/owr/2007/55, pp. 3209--3240: \url{https://doi.org/10.4171/owr/2007/55}. This is the sole supplied external reference.
\item Supplied literature triage, checked 2026-08-21: status ``open.'' It says both boundedness questions are posed unresolved in the source and reports no identified general resolution. This dated assessment derives from supplied OWR curation; time-sensitive status and cited claims require independent verification. It does not resolve the commuting-matrix defect in the extraction.
\end{itemize}
\end{document}
